\small
\textbf{Abstract.}
Retail forecasting models that see the price and promotion plan can explain
each forecast as a baseline, what would sell with nothing planned, plus the
lift each part of the plan adds. Merchants read that split as fact and plan on
it. We set out eight checks on such explanations, fixed in advance: properties
that hold by construction, sensitivity to randomised weights, stability across
seeds, price direction, agreement with what sold, agreement with LightGBM, a
plan read-back, and agreement with a measured cause. In generated store data
with known effects, a world model whose explanation is built in passes C1 and
C3 to C6, and its lifts match the truth
(rank agreement \kaClWmTruth) when plans are made at random. When merchants plan
into high-demand weeks, it still passes them, forecasts more accurately,
and overstates the total lift by a third (\kaCoWmBiasLo{} to \kaCoWmBiasHi{}
times) and the display's by \kaCoWmDispLo{} to \kaCoWmDispHi{} times; LightGBM
overstates them too, and the two agree. Asked about a lasting change in the
regular price, the world model answers with the temporary-cut response,
$\kaClWmCsevenLo$ to $\kaClWmCsevenHi$ against a true $-0.5$, and only the
measured cause finds it. The price-direction check catches every model blind to
price but also flags LightGBM; an exploratory lever-by-lever comparison catches
lifts given to the wrong lever; the randomisation check separates nothing. On
Dominick's we find \zEpisodes{} spells in which a store's regular prices sat
apart from its zone peers, most between April 1993 and April 1994. Over the
\zScored{} scored spells, prices \zSizeLo\% to \zSizeHi\% off their peers, the
pooled effect on category units is $\zSlope$ over eight weeks (95\% interval
$\zSlopeLo$ to $\zSlopeHi$) and $\zFullSlope$ over the full spell ($\zFullLo$ to
$\zFullHi$), in line with the elasticity near $-0.3$ the chain's experiments
implied. This is a noisy observational benchmark, not a ground truth: the
treated stores are inferred, prices move before the detected start, and a
placebo at false start dates gives $\zPlaceboSlope$ ($\zPlaceboLo$ to
$\zPlaceboHi$), close to the 8-week estimate.
We then audited the final tuned world model our companion paper reports for each
of its \tSets{} data sets (six Dominick's categories, Breakfast at the Frat, M5
and Favorita), \tWindows{} windows in all. Under the model's own plan for nothing
planned, the baseline and lift add up and the lift is zero on every window; the
lifts depend on what was learned (\tCTwoMean\% of windows) and the copies agree
(\tCThreeMean\%); and the plan can be read back from the state rolled under it
($R^2$ \tReadRolledLo{} to \tReadRolledHi). Price direction fails on every data set but
Frat: a deeper cut lowers the forecast for \tCFourDomLo\% to \tCFourDomHi\% of items on
Dominick's. Both the world model and LightGBM state about twice the measured
response to the zone price deviations (pooled slopes $\tSlopeWm$ and $\tSlopeLgb$
against $\tSlopeMeas$). Across the \tConfigsAll{} configurations the settings
searches tried, comparing those scored on the same data and windows, the number
of checks passed was not related to accuracy (pooled Spearman $\tRhoAll$, 95\%
interval $\tRhoAllLo$ to $\tRhoAllHi$). These
results hold for these models and data sets.

\vspace{4pt}
\textbf{Keywords:} retail forecasting; explanation; faithfulness; promotions;
price response; difference in differences; world models.
